@@M0@@
The result follows immediately from relation (1.11.3) and the induction hypothesis.

\sgaStableTarget{sga6.e06.prop.5-3}\textbf{Proposition 5.3.} \uline{With the notation and hypotheses of} 5.1, \uline{for} \uline{every} @@M1@@ \uline{there is a direct-sum decomposition of} @@M2@@-\uline{modules}:
@@M3@@
\uline{the filtration under consideration being the filtration of an augmented} @@M4@@-\uline{algebra defined} \uline{in} V 3.9.1 \uline{and} 3.10.

First note that @@M5@@ has augmentation zero. Moreover:
@@M6@@
Thus @@M7@@ for every @@M8@@. Consequently, @@M9@@ is a quotient of the @@M10@@-@@M11@@-algebra generated by a generator @@M12@@ subject to the relations @@M13@@ for every @@M14@@, and the ideal by which the quotient is taken is generated by relation (1.13.1). We can therefore apply V 4.12 and V 4.13.ii), and are reduced to showing that, for every @@M15@@:
@@M16@@

This is evident for @@M17@@; we continue by induction. Relation (1.13.1) gives:
@@M18@@
applying the induction hypothesis to the @@M19@@, we deduce the desired result, the sum being direct since @@M20@@ are independent over @@M21@@.

